\documentclass{article}
\usepackage{amssymb}
\usepackage{amsmath}
\usepackage{enumitem}

\begin{document}

  \section{Punto 1)}
    Demuestre que $D_n = \mathbb{Z}_2 \Join \mathbb{Z}_n$ encontrando el isomorfismo 
    $\phi$ correcto. Deduzca que el producto semidirecto de grupos abelianos no es abeliano.

    Vamos a deducir qué propiedades debería de cumplir $\phi$ para concluir el 
    isomorfismo entre $D_n$ y el producto semidirecto.

    Si es un homomorfismo, se debe seguir inmediatamente que $\phi(0) = I_{\mathbb{Z}_n}$ y 
    de la propiedad 2 enunciada, se debe seguir que $h k h^{-1} = \phi(h)(k)$. 

    Como los únicos elementos de $\mathbb{Z}_2$ son $0, 1$ y ya tenemos el valor de 
    $\phi$ en $0$, entonces la propiedad enunciada arriba se debe seguir para para 
    $1$, pero como queremos un isomorfismo con $D_n$, entonces podríamos decir que 
    $1$ sería el equivalente a la reflexión. 

    Así mismo, el $1'$ de $\mathbb{Z}_n$ sería el equivalente a la rotación de $D_n$. En 
    términos de parejas ordenadas, $s \sim (1, 0)$ y $r \sim (0, 1)$. 

    Como sabemos que $s r^j s = r^{-j}$, entonces se define
    \begin{equation*}
      \begin{aligned}
        \phi(1): \mathbb{Z}_n &\to \mathbb{Z}_n \\
        k &\to n - k 
      \end{aligned}
    \end{equation*}
    
    Básicamente, es el isomorfismo que toma un elemento y lo envía a su inverso.


    Debemos verificar que $\phi$ sea un homomorfismo, es decir, que $\phi(x + y) = \phi(x) \circ \phi(y)$.

    \begin{enumerate}[label=(\alph*)]
      \item Si $x \neq y$: Entonces alguno de los dos es el $0$ de $\mathbb{Z}_2$, luego, 
            $\phi(x + y) = \phi(y)$ o $\phi(x + y) = \phi(x)$, en cualquier caso, se cumple,
            pues el contrarío sería la función identidad y se concluiría que 
            $\phi(x + y) = \phi(x) \circ \phi(y)$.

      \item Si $x = y$, entonces $x + y = 0$. Si $x = 0$, es inmediato. Si $x = 1$, entonces 
            bastaría probar que $\phi(1) \circ \phi(1) = I_{\mathbb{Z}_n}$. 
            $(\phi(1) \circ \phi(1)) (k) = \phi(1)(\phi(1)(k)) = \phi(1)(n-k) = k$, concluyendo así que 
            $\phi(x + y) = \phi(x) \circ \phi(y)$.
    \end{enumerate}

    Con todo esto, vamos a ver ahora si que $D_n \equiv \mathbb{Z}_2 \Join \mathbb{Z}_n$.

    Definimos la siguiente función:
    \begin{equation*}
      \begin{aligned}
        \Psi: D_n &\to \mathbb{Z}_2 \Join \mathbb{Z}_n \\
        s^i \cdot r^j &\to (i, \phi(i)(j))
      \end{aligned}
    \end{equation*}

    Verifiquemos que esta es un isomorfismo. Primero, verifiquemos que $\Psi(s^i r^j) = \Psi(s^i) \Psi(r^j)$.

    \begin{equation*}
      \begin{aligned}
        \Psi(s^i r^j) &= (i, \phi(i)(j)) \\
        \Psi(s^i) \Psi(r^j) &= (i, \phi(i)(0)) \cdot (0, \phi(0)(j)) \\
        &= (i, 0) \cdot (0, j) \\
        &= (i + 0, 0 + \phi(i)(j)) \\
        &= (i, \phi(i)(j))
      \end{aligned}
    \end{equation*}

    Y como todo elemento de $D_n$ tiene esta forma, concluimos que es un homomorfismo sobreyectivo (cualquier 
    pareja ordenada $(i, k)$ tiene pre-imagen $s^i r^{(\phi(i)^{-1}(k)}$, pues $\Psi(s^i r^{(\phi(i)^{-1}(k)}) = 
    (i, \phi(i)((\phi(i)^{-1}(k)) = (i, k)$.


    Ahora, si $\Psi(s^i r^j) = (0, 0)$, entonces $(i, \phi(i)(j)) = (0, 0)$, luego $i = 0$, entonces 
    $j = \phi(i)(j) = 0$, entonces es un homomorfismo biyectivo, luego es un isomorfismo.



    Finalmente, sabemos que tanto $\mathbb{Z}_2$ y $\mathbb{Z}_n$ son grupos abelianos, pero $D_n$ no 
    es un grupo abeliano, concluyendo así que el producto semidirecto de grupos abelianos no es, necesariamente, 
    abeliano.
\end{document}
